%% kapa: exam=1 round=29 year=2018 period=1 form=A subject=economics subject_ko=경제학원론 seq=35 number=75
%% source: Q-Net 감정평가사 기출문제 (cst00309) artlSeq=5202808 file=extracted/1008/5202808/2018년 제29회 감정평가사 1차시험 기출문제/2018년 제29회 감정평가사 1차시험 기출문제/1교시 A형.pdf
%% conversion: kordoc 4.24.51 -> tools/gen1.py (auto); review=proofread-sonnet; flags=none
\begin{question}{75}
모든 시장이 완전경쟁 상태인 경제에서 총생산함수는 $Y=AL^{2/3}K^{1/3}$ 이다. 매년 $L$, $K$, $A$가 각각 3 \%씩 증가하는 경제에 관한 설명으로 옳은 것을 모두 고른 것은? (단, $Y$는 국내총생산, $L$은 노동량, $K$는 자본량, $A$는 상수이다.)
\begin{framedtext}
ㄱ. 총생산함수는 규모 수익 불변이다. \par
ㄴ. 노동소득분배율은 $2/3$이다. \par
ㄷ. 경제성장률은 6 \%이다. \par
\end{framedtext}
\begin{choices}
  \item ㄱ
  \item ㄴ
  \item ㄱ, ㄴ
  \item ㄴ, ㄷ
  \item ㄱ, ㄴ, ㄷ
\end{choices}
\end{question}
\answer{5}
